\answer{ex:05:1}{$\ell\approx17\um$; $\approx100$~діб: розсилку по мозку робить розгалуження проводу, дифузія лише розмиває кожне місце викиду на $\sim20\um$.}
\answer{ex:05:2}{$r=ab/(1+G)$, $S=1/(1+G)$ точно; $+1.7\,\%$ замість $+20\,\%$.}
\answer{ex:05:3}{$T^*=1.21\,\text{с}$ при $G=2$ і $0.19\,\text{с}$ при $G=9$: реально $G\sim2$--$4$, пригнічення в $3$--$5$ разів.}
\answer{ex:05:4}{$\mathrm{CV}=1/\sqrt{2\lambda\tau_c}\approx9\cdot10^{-4}$ за сумарної частоти $\lambda$; одному нейрону потрібно $\tau_c=12.5\,\text{с}$.}
\answer{ex:05:5}{$c\approx0.40$ в обох випадках; $\mathrm{CV}=0.050$ проти $0.158$, у $\sqrt{10}$ разів. Частоти окремих нейронів лишаються пропорційними $a_i$: регулюється лише сума.}
\answer{ex:05:6}{$N_1=\overline{A\vee B}$, $N_2=\overline{A\vee N_1}$, $N_3=\overline{B\vee N_1}$, $N_4=\overline{N_2\vee N_3}$, XOR $=N_5=\overline{N_4}$; кожне перемикання займає $\tau_c\ln2$, шлях із $4$ вентилів --- $2.8\,\text{с}$ на один XOR.}
\answer{ex:05:7}{(а)~$\tau_\gamma=6/\ln3\approx5.5\,\text{с}$. (б)~$\bar D<4\,\text{с}$ проти $D<2.2\,\text{с}$: непередбачуваність затримки робить терплячішим.}
\answer{ex:05:8}{$k_2=1/\tau_d=10\,\text{с}^{-1}$, $k_1=1/\tau_d-1/\tau_\gamma=9.5\,\text{с}^{-1}$; $\tau_\gamma=1/(k_2-mk_1)$: подвоєння при $+2.6\,\%$ (а при $+5.3\,\%$ горизонт нескінченний).}
